\documentclass[aspectratio=169]{beamer}
\usetheme{Madrid}
\usecolortheme{default}

\usepackage[utf8]{inputenc}
\usepackage[T5]{fontenc}
\usepackage[vietnamese]{babel}
\usepackage{graphicx}
\usepackage{amsmath,amssymb}
\usepackage{booktabs}
\usepackage{tikz}
\usetikzlibrary{shapes,arrows,positioning}
\usepackage{listings}
\usepackage{xcolor}
\usepackage{multicol}

\definecolor{codegreen}{rgb}{0,0.6,0}
\definecolor{codegray}{rgb}{0.5,0.5,0.5}
\definecolor{codepurple}{rgb}{0.58,0,0.82}
\definecolor{backcolour}{rgb}{0.95,0.95,0.92}

\lstdefinestyle{mystyle}{
    backgroundcolor=\color{backcolour},
    commentstyle=\color{codegreen},
    keywordstyle=\color{magenta},
    numberstyle=\tiny\color{codegray},
    stringstyle=\color{codepurple},
    basicstyle=\ttfamily\footnotesize,
    breaklines=true,
    frame=single
}
\lstset{style=mystyle}

\title[Trích xuất Hóa đơn VN]{Trích xuất Thông tin Hóa đơn Tiếng Việt từ Ảnh}
\subtitle{Pipeline PaddleOCR -- VietOCR -- LayoutLMv3}
\author{Đồ án Môn học}
\institute{Học viện Công nghệ Bưu chính Viễn thông \\ Khoa Công nghệ Thông tin 2}
\date{Tháng 6, 2026}

\begin{document}

%---------------------------------------------------------
% Slide 1: Title
%---------------------------------------------------------
\begin{frame}
    \titlepage
\end{frame}

%---------------------------------------------------------
% Slide 2: Outline
%---------------------------------------------------------
\begin{frame}{Nội dung trình bày}
    \tableofcontents
\end{frame}

%---------------------------------------------------------
% Section 1: MỞ ĐẦU
%---------------------------------------------------------
\section{Mở đầu}

% Slide 3
\begin{frame}{Bối cảnh bài toán}
    \begin{itemize}
        \item Mỗi ngày, hàng triệu giao dịch thương mại tại Việt Nam tạo ra hàng triệu tờ hóa đơn giấy
        \item Kế toán viên gõ tay từng con số $\rightarrow$ sai sót nhập liệu thường xuyên
        \item Máy tính đã có thể đọc chữ in từ lâu, nhưng đọc hóa đơn thực tế vẫn là thách thức:
        \begin{itemize}
            \item Font chữ tùy ý, góc chụp nghiêng
            \item Ánh sáng không đều
            \item Tiếng Việt với đầy đủ dấu thanh
        \end{itemize}
        \item Cần biến \textbf{ảnh hóa đơn} thành \textbf{dữ liệu có cấu trúc}
    \end{itemize}
\end{frame}

% Slide 4
\begin{frame}{Các tình huống thực tế}
    \textbf{Kịch bản A -- Kế toán doanh nghiệp nhỏ:}
    \begin{itemize}
        \item 300 hóa đơn công tác phí, mất 2--3 ngày nhập thủ công
    \end{itemize}
    \vspace{0.3cm}
    \textbf{Kịch bản B -- Ứng dụng chi tiêu cá nhân:}
    \begin{itemize}
        \item Chụp ảnh hóa đơn $\rightarrow$ tự động điền ngày, siêu thị, tổng tiền
    \end{itemize}
    \vspace{0.3cm}
    \textbf{Kịch bản C -- Kiểm toán \& tuân thủ thuế:}
    \begin{itemize}
        \item Xác minh hàng nghìn hóa đơn VAT tự động
    \end{itemize}
    \vspace{0.3cm}
    $\Rightarrow$ Nhu cầu chung: \textbf{Document Information Extraction (DIE)}
\end{frame}

%---------------------------------------------------------
% Section 2: TỔNG QUAN
%---------------------------------------------------------
\section{Tổng quan bài toán}

% Slide 5
\begin{frame}{Mô hình hóa bài toán -- Ba tầng xử lý}
    \begin{center}
    \begin{tikzpicture}[
        box/.style={rectangle, draw, rounded corners, fill=blue!15, text width=11cm, minimum height=1.2cm, align=center, font=\small},
        arrow/.style={->, thick, >=stealth}
    ]
        \node[box] (t3) at (0,3.5) {\textbf{Tầng 3:} Trích xuất thực thể có cấu trúc (KIE) -- LayoutLMv3};
        \node[box] (t2) at (0,1.8) {\textbf{Tầng 2:} Nhận dạng ký tự (OCR Recognition) -- VietOCR};
        \node[box] (t1) at (0,0) {\textbf{Tầng 1:} Phát hiện vùng chữ (Text Detection) -- PaddleOCR};
        \draw[arrow] (t1) -- (t2);
        \draw[arrow] (t2) -- (t3);
    \end{tikzpicture}
    \end{center}
\end{frame}

% Slide 6
\begin{frame}{Đầu vào \& Đầu ra}
    \textbf{Đầu vào:} Ảnh JPEG/PNG hóa đơn chụp bằng điện thoại (nghiêng, nhòe, ánh sáng không đều)
    \vspace{0.4cm}

    \textbf{Đầu ra:} Cấu trúc JSON có ngữ nghĩa
    \begin{lstlisting}[language=Python]
{
  "store_name": "Circle K Nguyen Hue",
  "date": "2024-05-27",
  "total": 87000,
  "items": [
    {"name": "Ca phe sua da", "qty": 2, "price": 29000},
    {"name": "Banh mi que", "qty": 1, "price": 15000}
  ]
}
    \end{lstlisting}
\end{frame}

% Slide 7
\begin{frame}{Pipeline xử lý}
    \begin{center}
    \begin{tikzpicture}[
        node distance=1.5cm,
        process/.style={rectangle, draw, rounded corners, fill=blue!10, minimum height=1cm, minimum width=2.5cm, align=center, font=\footnotesize},
        data/.style={rectangle, draw, rounded corners, fill=green!10, minimum height=0.8cm, minimum width=2.5cm, align=center, font=\footnotesize},
        arrow/.style={->, thick, >=stealth}
    ]
        \node[data] (img) {Ảnh hóa đơn};
        \node[process, right=of img] (det) {PaddleOCR\\Detection};
        \node[data, right=of det] (bbox) {Bounding\\Boxes};
        \node[process, right=of bbox] (rec) {VietOCR\\Recognition};
        \node[data, right=of rec] (text) {Text + BBox};
        \node[process, below=1.5cm of rec] (kie) {LayoutLMv3\\KIE};
        \node[data, left=of kie] (json) {JSON\\Cấu trúc};

        \draw[arrow] (img) -- (det);
        \draw[arrow] (det) -- (bbox);
        \draw[arrow] (bbox) -- (rec);
        \draw[arrow] (rec) -- (text);
        \draw[arrow] (text) -- (kie);
        \draw[arrow] (kie) -- (json);
    \end{tikzpicture}
    \end{center}
\end{frame}

%---------------------------------------------------------
% Section 3: PHÂN LOẠI & TIẾN HÓA
%---------------------------------------------------------
\section{Phân loại và tiến hóa}

% Slide 8
\begin{frame}{Text Detection -- Từ sliding window đến anchor-free}
    \begin{columns}[T]
        \begin{column}{0.48\textwidth}
            \textbf{Thế hệ 1 (trước 2015):}
            \begin{itemize}
                \item MSER, SWT
                \item Phân tích pixel theo màu sắc, độ rộng nét bút
                \item Chỉ hoạt động với ảnh sạch, chữ đứng
            \end{itemize}
            \vspace{0.3cm}
            \textbf{Thế hệ 2 (2015--2017):}
            \begin{itemize}
                \item CTPN, TextBoxes++, SegLink
                \item CNN-based với Anchor Boxes
                \item Khó phát hiện văn bản cong
            \end{itemize}
        \end{column}
        \begin{column}{0.48\textwidth}
            \textbf{Thế hệ 3 (2018--nay):}
            \begin{itemize}
                \item \textbf{DB (Differentiable Binarization)}
                \item Xương sống của PaddleOCR
                \item Dự đoán probability map per pixel
                \item Ngưỡng có thể học được (differentiable threshold)
                \item Tối ưu cho CPU/mobile
            \end{itemize}
        \end{column}
    \end{columns}
\end{frame}

% Slide 9
\begin{frame}{Text Recognition -- Từ template matching đến Attention}
    \begin{columns}[T]
        \begin{column}{0.48\textwidth}
            \textbf{Thế hệ 1 (trước 2014):}
            \begin{itemize}
                \item Tesseract v3, template matching
                \item Tách từng ký tự, so sánh template
                \item Yêu cầu font rõ ràng
            \end{itemize}
            \vspace{0.3cm}
            \textbf{Thế hệ 2 (2015--2018):}
            \begin{itemize}
                \item \textbf{CRNN + CTC}
                \item CNN + BiLSTM + CTC loss
                \item Nền tảng OCR hiện đại
            \end{itemize}
        \end{column}
        \begin{column}{0.48\textwidth}
            \textbf{Thế hệ 3 (2018--nay):}
            \begin{itemize}
                \item \textbf{VietOCR -- Seq2Seq + Attention}
                \item VGG/ResNet encoder + Transformer decoder
                \item Xử lý dấu thanh tiếng Việt tốt hơn
                \item Pretrain trên 5M+ ảnh text tiếng Việt
                \item 60+ font chữ thông dụng
            \end{itemize}
        \end{column}
    \end{columns}
\end{frame}

% Slide 10
\begin{frame}{Key Information Extraction -- Từ rule-based đến Multi-modal}
    \begin{columns}[T]
        \begin{column}{0.48\textwidth}
            \textbf{Thế hệ 1 -- Rule-based:}
            \begin{itemize}
                \item Regex, heuristic
                \item Nhanh, không cần training
                \item Vỡ ngay khi gặp format mới
            \end{itemize}
            \vspace{0.3cm}
            \textbf{Thế hệ 2 -- BERT NER:}
            \begin{itemize}
                \item Token classification
                \item Bỏ qua thông tin vị trí không gian
            \end{itemize}
        \end{column}
        \begin{column}{0.48\textwidth}
            \textbf{Thế hệ 3 -- LayoutLMv3:}
            \begin{itemize}
                \item Microsoft, 2022
                \item Kết hợp 3 nguồn: text OCR + tọa độ bbox + patch ảnh
                \item Unified self-attention qua cả 3 modality
                \item Pre-training: MLM + MIM + Text-Image Alignment
                \item \textbf{Đột phá} cho Document AI
            \end{itemize}
        \end{column}
    \end{columns}
\end{frame}

%---------------------------------------------------------
% Section 4: TRỰC GIÁC
%---------------------------------------------------------
\section{Trực giác và tư duy}

% Slide 11
\begin{frame}{Con người đọc hóa đơn như thế nào?}
    \begin{enumerate}
        \item \textbf{Vị trí không gian mang ngữ nghĩa:}
        \begin{itemize}
            \item ``Tổng cộng'' $\rightarrow$ góc dưới phải, font lớn, in đậm
            \item Tên cửa hàng $\rightarrow$ trên cùng, canh giữa
        \end{itemize}
        \item \textbf{Căn lề là thông tin:}
        \begin{itemize}
            \item Giá tiền căn phải, tên mặt hàng căn trái
        \end{itemize}
        \item \textbf{Gần nhau $\rightarrow$ liên quan:}
        \begin{itemize}
            \item ``2 x 29,000'' và ``58,000'' cùng dòng $\rightarrow$ cùng mặt hàng
        \end{itemize}
        \item \textbf{Pattern lặp lại trong vùng bảng:}
        \begin{itemize}
            \item 5 dòng có cấu trúc [tên $|$ số $|$ số] $\rightarrow$ bảng mặt hàng
        \end{itemize}
    \end{enumerate}
\end{frame}

% Slide 12
\begin{frame}{Mô hình mô phỏng trực giác con người}
    \begin{center}
    \begin{tabular}{p{3.5cm}p{8cm}}
        \toprule
        \textbf{Trực giác} & \textbf{Mô hình mô phỏng} \\
        \midrule
        Vùng chữ có cấu trúc pixel đặc biệt & PaddleOCR Detection: học probability map \\
        \addlinespace
        Đọc ký tự cần ngữ cảnh xung quanh & VietOCR: attention ``nhìn lại'' toàn bộ ảnh \\
        \addlinespace
        Vị trí + nội dung + hình ảnh cùng quyết định & LayoutLMv3: fuse 3 luồng trong transformer \\
        \bottomrule
    \end{tabular}
    \end{center}
    \vspace{0.5cm}
    \textbf{Giả định ngầm:} Mỗi mô hình đều có giả định riêng về dữ liệu input $\rightarrow$ khi giả định vi phạm, mô hình suy giảm chất lượng.
\end{frame}

%---------------------------------------------------------
% Section 5: TOÁN HỌC HÓA
%---------------------------------------------------------
\section{Toán học hóa}

% Slide 13
\begin{frame}{PaddleOCR Detection -- DB (Differentiable Binarization)}
    \textbf{Phép toán chủ đạo:} Segmentation + Adaptive Thresholding
    \vspace{0.3cm}

    Mô hình học 2 map song song từ feature map $F$:
    \begin{align*}
        P &= \sigma(\text{Conv}(F)) && \leftarrow \text{Probability map} \\
        T &= \sigma(\text{Conv}'(F)) && \leftarrow \text{Threshold map} \\
        \hat{B} &= \sigma\left(\frac{P - T}{k}\right) && \leftarrow \text{Binary map (differentiable)}
    \end{align*}
    \vspace{0.2cm}

    \textbf{Bí quyết:} $\sigma\left(\frac{P-T}{k}\right)$ với $k=50$ xấp xỉ hàm step nhưng vẫn có gradient để backprop.

    \textbf{Loss:} $L = L_{\text{bce}}(P, G_p) + \alpha \cdot L_{\text{bce}}(\hat{B}, G_b) + \beta \cdot L_1(T, G_t)$

    \textbf{Tham số:} ResNet-18/50 + FPN, khoảng 10--27M tham số.
\end{frame}

% Slide 14
\begin{frame}{VietOCR -- Seq2Seq với Attention}
    \textbf{Phép toán chủ đạo:} Cross-Attention trong Transformer decoder
    \vspace{0.3cm}

    \begin{align*}
        F &= \text{CNN\_Encoder}(\text{ảnh}) \quad \text{shape: } (H/4, W/4, C) \\
        F_{\text{flat}} &= \text{reshape}(F) + \text{positional\_embed} \quad \text{shape: } (S, C) \\
        h_t &= \text{TransformerDecoder}(h_{t-1}, F_{\text{flat}}) \\
        p(y_t) &= \text{Softmax}(\text{Linear}(h_t))
    \end{align*}
    \vspace{0.2cm}

    \textbf{Attention:} $\alpha_{t,s} = \text{softmax}\left(\frac{Q_t \cdot K_s}{\sqrt{d_k}}\right)$ -- cho phép ``nhìn'' vào vùng feature tương ứng.

    \textbf{Đặc biệt tiếng Việt:} Khi sinh ``ề'', decoder attend mạnh vào vùng có dấu phụ kép (mũ + huyền).

    \textbf{Vocab:} 130--150 ký tự Unicode tiếng Việt + dấu câu + số.
\end{frame}

% Slide 15
\begin{frame}{LayoutLMv3 -- Multi-modal Transformer}
    \textbf{Phép toán chủ đạo:} Unified Self-Attention qua cả 3 modality
    \vspace{0.3cm}

    \textbf{Input token $i$:}
    \[e_i = E_{\text{text}}(\text{token}_i) + E_{\text{1D}}(\text{pos}_i) + E_{\text{2D}}(x_1, y_1, x_2, y_2, w, h)\]

    \textbf{Image patch $j$:}
    \[v_j = \text{LinearProj}(\text{patch}_j) + E_{\text{img\_pos}}(j)\]

    \textbf{Concatenate $\rightarrow$ Transformer $\rightarrow$ Contextualized representations}
    \vspace{0.3cm}

    \textbf{Pre-training objectives:}
    \begin{itemize}
        \item Masked Language Modeling (MLM) trên text
        \item Masked Image Modeling (MIM) trên ảnh
        \item Text-Image Alignment
    \end{itemize}
\end{frame}

%---------------------------------------------------------
% Section 6: KHOẢNG CÁCH
%---------------------------------------------------------
\section{Khoảng cách giữa lý thuyết và thực tế}

% Slide 16
\begin{frame}{Hàm mất mát và lý do lựa chọn}
    \begin{center}
    \begin{tabular}{p{3cm}p{4cm}p{4.5cm}}
        \toprule
        \textbf{Tầng} & \textbf{Loss Function} & \textbf{Tại sao?} \\
        \midrule
        Detection & BCE + L1 & Gradient ổn định, robust với outlier \\
        \addlinespace
        Recognition & Cross-Entropy chuỗi & Smooth, vi phân được \\
        \addlinespace
        KIE & Cross-Entropy/token & Tính được tại từng token \\
        \bottomrule
    \end{tabular}
    \end{center}
    \vspace{0.4cm}

    \textbf{Tại sao không dùng Metric làm Loss?}
    \begin{itemize}
        \item Field Exact Match: không có gradient
        \item CER (edit distance): không liên tục
        \item Entity F1: phụ thuộc ngưỡng confidence
    \end{itemize}
    $\Rightarrow$ Dùng CE làm Loss, F1/CER/Exact Match chỉ để đánh giá.
\end{frame}

% Slide 17
\begin{frame}{Hệ quả thực tế}
    \begin{alertblock}{Điểm đau lớn nhất}
        Mô hình có thể đạt \textbf{CER rất thấp (98\% ký tự đúng)} nhưng \textbf{Field Exact Match chỉ 70\%} -- vì chỉ cần sai 1 ký tự trong ``87.000'' thành ``87.500'' là cả trường total sai hoàn toàn.
    \end{alertblock}
    \vspace{0.3cm}

    \textbf{Khả năng mở rộng Loss:}
    \begin{itemize}
        \item Thêm class mới: DISCOUNT, TAX\_CODE
        \item Weighted loss: ưu tiên TOTAL, DATE
        \item Consistency constraint: $\sum$(subtotals) $\approx$ total
    \end{itemize}
\end{frame}

%---------------------------------------------------------
% Section 7: THÁCH THỨC
%---------------------------------------------------------
\section{Thách thức và rào cản}

% Slide 18
\begin{frame}{Trở ngại đã giải quyết \& còn tồn đọng}
    \textbf{Đã giải quyết:}
    \begin{itemize}
        \item Văn bản đa hướng, cong $\rightarrow$ DB, ABCNet (bezier curves)
        \item Font chữ đa dạng tiếng Việt $\rightarrow$ VietOCR pretrained + augmentation
        \item Hiểu layout cơ bản $\rightarrow$ 2D positional embedding (LayoutLM)
    \end{itemize}
    \vspace{0.3cm}

    \textbf{Còn tồn đọng:}
    \begin{itemize}
        \item[$\times$] Ảnh chất lượng thấp (run, nhòe, phản chiếu) -- CER tụt 85--90\%
        \item[$\times$] Hóa đơn in nhiệt bị phai -- contrast rất thấp
        \item[$\times$] Zero-shot generalization -- cần fine-tune cho mỗi template mới
        \item[$\times$] Dữ liệu tiếng Việt khan hiếm, hạ tầng GPU đắt
        \item[$\times$] Relation Extraction phức tạp (20+ mặt hàng)
    \end{itemize}
\end{frame}

% Slide 19
\begin{frame}{Hướng đi tương lai}
    \begin{enumerate}
        \item \textbf{End-to-end Document Understanding:}
        \begin{itemize}
            \item Donut, Nougat -- bỏ qua OCR, ảnh $\rightarrow$ JSON trực tiếp
        \end{itemize}
        \item \textbf{Large Vision-Language Models (LVLMs):}
        \begin{itemize}
            \item GPT-4V, Gemini -- zero-shot, không cần fine-tune
        \end{itemize}
        \item \textbf{Self-supervised pretraining trên tài liệu Việt:}
        \begin{itemize}
            \item Corpus hóa đơn/biểu mẫu tiếng Việt quy mô lớn
        \end{itemize}
        \item \textbf{Synthetic data generation:}
        \begin{itemize}
            \item Tạo hóa đơn giả + domain randomization
        \end{itemize}
    \end{enumerate}
\end{frame}

%---------------------------------------------------------
% Section 8: KẾT LUẬN
%---------------------------------------------------------
\section{Kết luận}

% Slide 20
\begin{frame}{Kết luận}
    \begin{itemize}
        \item Bài toán trích xuất hóa đơn tiếng Việt đòi hỏi giải quyết đồng thời 3 thách thức: \textbf{phát hiện vùng chữ}, \textbf{nhận dạng ký tự}, \textbf{hiểu cấu trúc ngữ nghĩa}
        \vspace{0.2cm}
        \item Pipeline \textbf{PaddleOCR $\rightarrow$ VietOCR $\rightarrow$ LayoutLMv3} là lựa chọn hợp lý hiện tại với pretrained weights chất lượng tốt
        \vspace{0.2cm}
        \item Nhược điểm cố hữu: \textbf{lỗi tích lũy qua từng tầng}
        \vspace{0.2cm}
        \item Ranh giới giữa OCR+NLP pipeline và end-to-end LVLM đang dần mờ đi
    \end{itemize}
\end{frame}

%---------------------------------------------------------
% Slide 21: Cảm ơn
%---------------------------------------------------------
\begin{frame}{}
    \begin{center}
        \vspace{1.5cm}
        {\Huge\bfseries Cảm ơn Thầy và Các bạn!}
        \vspace{1cm}

        {\Large Xin cảm ơn!}
        \vspace{0.5cm}

        {\normalsize Học viện Công nghệ Bưu chính Viễn thông \\ Khoa Công nghệ Thông tin 2}
    \end{center}
\end{frame}

\end{document}
